


\documentclass[11pt]{article}
\usepackage{amsgen,amsmath,amstext,amsbsy,amsopn,amsfonts,amsthm,amscd}
%\usepackage{bbm}
\usepackage[dvips]{graphicx}
\usepackage{enumerate}
\setlength{\textwidth}{16.0cm} \setlength{\textheight}{25cm}
\voffset=-2.50cm \hoffset=-2.00cm \setlength{\topmargin}{0.4cm}

%%%%%%%%%%%%%%%%%%%%%%%%%%%%%%%%%%%%%%%%%%%%%%%%%%%%%%%%%%%%%%

\newtheorem{lem1}{Lema}%[chapter]
\newtheorem{th1}{Teorema}%[chapter]
\newtheorem{prop1}{Proposici\'{o}n}%[chapter]
\newtheorem{defi1}{Definici\'{o}n}%[chapter]
\newtheorem{coro1}{Corolario}%[chapter]
\newtheorem{ob1}{Observaci\'{o}n}%[chapter]
\newtheorem{obs1}{Observaci\'{o}n}%[chapter]
\newtheorem{nota1}{Nota}%[chapter]
\newtheorem{eje1}{Ejemplo}%[chapter]
\newtheorem{ejes1}{Ejemplo}%[chapter]
\newtheorem{ejerc1}{Ejercicio}%[chapter]
\newtheorem{ejercs1}{Ejercicio}%[chapter]
\newenvironment{prop}{\begin{prop1}\rm}{\end{prop1}}
\newenvironment{defi}{\begin{defi1}}{\end{defi1}}
\newenvironment{coro}{\begin{coro1}}{\end{coro1}}
\newenvironment{th2}{\begin{th1}}{\end{th1}}
\newenvironment{nota}{\begin{nota1}\rm}{\end{nota1}}
\newenvironment{lem}{\begin{lem1}}{\end{lem1}}
\newenvironment{ob}{\begin{ob1}\rm}{\end{ob1}}
\newenvironment{obs}{\begin{obs1}\rm}{\end{obs1}}
\newenvironment{eje}{\begin{eje1}\rm}{\end{eje1}}
\newenvironment{ejes}{\begin{ejes1}\rm}{\end{ejes1}}
\newenvironment{ejerc}{\begin{ejerc1}\rm}{\end{ejerc1}}
\newenvironment{ejercs}{\begin{ejercs1}\rm}{\end{ejercs1}}
%%%%%%%%%%%%%%%%%%%%%%%%%%%%%%%%%%%%%%%%%%%%%%%%%%%%%%
%\DeclareSymbolFont{cucu}{U}{bbmss}{m}{n}
\DeclareMathOperator{\CC}{\mathbb C}%{\mathalpha}{cucu}{67}
\DeclareMathOperator{\NN}{\mathbb N}%{\mathalpha}{cucu}{78}
\DeclareMathOperator{\RR}{\mathbb R}%{\mathalpha}{cucu}{82}
\DeclareMathOperator{\QQ}{\mathbb Q}%{\mathalpha}{cucu}{81}
\DeclareMathOperator{\ZZ}{\mathbb Z}%{\mathalpha}{cucu}{90}
%%%%%%%%%%%%%%%%%%%%%%%%%%%%%%%%%%%%%%%%%%%%%%%%%%%%%%
\DeclareMathOperator{\sen}{sen}
\DeclareMathOperator{\arcsen}{arcsen}
\DeclareMathOperator{\senh}{senh} \DeclareMathOperator{\fr}{Frac}
%%%%%%%%%%%%%%%%%%%%%%%%%%%%%%%%%%%%%%%%%%%%%%%%%%%%
\newcommand{\ran}[1]{{\em Im\/}(#1)}
\newcommand{\ca}[1]{{\em card\/}(#1)}
\newcommand{\sop}[1]{{\em supp\/}(#1)}
\newcommand{\tra}[1]{{\em traza\/}(#1)}
\newcommand{\dis}[1]{d (#1)}
\newcommand{\inte}[1]{{\em Int\/}(#1)}
\newcommand{\conv}[1]{{\em conv\/}(#1)}
\newcommand{\too}{\longrightarrow}
\newcommand{\oo}{\overline}
\newcommand{\cq}{\begin{flushright} $\Box$ \end{flushright}}
\newcommand{\cqd}{\, \, \rule{7pt}{7pt}}
\def\limsup{\mathop{\overline{\rm lim}}}
\def\liminf{\mathop{\underline{\rm lim}}}
\renewcommand{\theenumii}{\arabic{enumii}}
\renewcommand{\labelenumii}{\theenumi.\theenumii.}
\renewcommand{\theenumiii}{\arabic{enumiii}}
\renewcommand{\labelenumiii}{\theenumi.\theenumii.\theenumiii}
%%%%%%%%%%%%%%%%%%%%%%%%%%%%%%%%%%%%%%%%%%%%%%%%%%%%%%%%%%%%%%%%%
\begin{document}

\section{SEGUNDO VALOR SINGULAR}



\noindent $\bullet$ {\bf MOTIVACIÓN:  MUY INTERESANTE:}

\noindent Estudio de la convergencia del \underline{segundo} mayor   de los valores singulares de $\mathbf{D}_n$, cuyo comportamiento, en los experimentos realizados, no parece ser  el mismo que el del primer autovalor!!!!. y que puede dar una idea más aproximada del soporte: de hecho parece converger al mayor valor absoluto de los elementos del soporte.

  {\bf Problemas relacionados}
\begin{enumerate}
  \item  Interesante: puede ocurrir que al tomar la primera derivada vamos al punto límite del soporte, al tomar la segunda derivada son los terceros máximos valores singulares los que van a dicho valor ????
  \item Cuál es la situación para los valores singulares cuando consideramos el caso standard de medidas (no sobolev) y caso real.
        ????
  \item Idea curiosa: cuando en un espacio con peso en $0$ consideramos el operador multiplicación restringido al subespacio vectorial de polinomios $[z^{2}p(z), p(z)\in \mathbf{P}[z]]$ se obtiene que la norma del operador multiplicación con las normas sobolev coincide con las normas del operador multiplicación con la norma asociada a la medida.

\end{enumerate}


\subsection{Una introducción a los valores singulares de una matriz}

\noindent Intoducimos los valores singulares de una matriz cualquiera finita:

\noindent Sea $\mathbf{A}$ una matriz $n\times n$ de números complejos,  la matriz $\mathbf{A}^{*}\mathbf{A}$  ( $\mathbf{A}^{t}\mathbf{A}$ en el caso de matrices con entradas reales) es una matriz hermitiana y semidefinida positiva. Por lo tanto, tiene $n$ autovalores mayores o iguales que cero contados con su multiplicidad
$$
  0 \leq \lambda_{n,n}\leq \dots \leq \lambda_{2,n} \leq \lambda_{1,n}
$$


\noindent Los valores singulares de $\mathbf{A}$ son:
$$
  \sigma_{k,n}= \sqrt{\lambda_{k,n}}
$$
\noindent siendo:
$$
  0 \leq \sigma_{n,n} \leq \dots \sigma_{2,n} \leq \sigma_{1,n}
$$

\noindent Utilizando el criterio de Rayleigh y el Principio de min-max de  Courant-Fischer  (BUSCAR BIBLIOGRAFÍA) se tiene que estos valores singulares tienen relación con la norma de aplicaciones lineales y continuas cuando consideramos la norma euclídea en los espacios $\CC^{n}$. De hecho,
\begin{enumerate}
  \item
        $$
          \sigma_{1,n}^2 =
          \max \{ <\mathbf{A}^{t}\mathbf{A}x,x>: \Vert x \Vert =1 \}
          =
          = \max \{\Vert \mathbf{A} x \Vert : \Vert x \Vert =1\}=
          \Vert \mathbf{A} \Vert^2
        $$
  \item Principio min-max: el segundo mayor valor singular de una matriz es el mínimo de las normas de las restricciones de la aplicación $\mathbf{A}$ a los espacios de codimnesión $1$.
        $$
          \sigma_{2,n}^2 = \min_{V: cod(V)=1} \max \{\Vert \mathbf{A}x \Vert , x\in V, \Vert x \Vert =1\}= \min_{V: cod(V)=1} \Vert \mathbf{A}/V \Vert
        $$

        \noindent siendo $\mathbf{A}/V$ la restricción de la aplicación lineal $\mathbf{A}$ al subespacio $V$.

\end{enumerate}







\subsection{Convergencia del mayor valor singular de $\mathbf{D}_n$.}





\noindent Dada la matriz $\mathbf{D}_n^{t}\mathbf{D}_n$ que es simétrica y semidefinida positiva consideramos el segundo mayor valor singular, es decir, si ordenamos los valores singulares de esta forma:

$$
  0\leq \sigma_{n,n} \leq \dots \leq \sigma_{2,n} \leq \sigma_{1,n}
$$



\noindent son los valores singulares de $\mathbf{D}_n$,

\bigskip

\noindent Nótese que $\Vert \mathbf{D}_n \Vert= \sigma_{1,n}$, ya que para una matriz simétrica y semidefinida positiva el mayor autovalor es la norma.

\bigskip Utilizando resultados de Análisis Funcional (BUSCAR REFERENCIA) se tiene que:

\begin{prop} Sea $\mathcal{D}$ una aplicación lineal continua en $\ell_2$ y sean $\mathbf{D}_n$ las secciones de orden $n$, entonces
  $$
    \Vert \mathcal{D} \Vert = \lim_{n\to \infty} \Vert \mathbf{D}_n\Vert
  $$

\end{prop}

\noindent Utilizando este resultado se tiene que:
$$
  \lim_{n\to \infty} \sigma_{1,n}(\mathbf{D}_n)= \lim_{n \to \infty} \Vert \mathbf{D}_n\Vert= \Vert \mathcal{D}\Vert
$$

\noindent que siempre existe pudiendo ser $\infty$ en el caso no acotado.

\subsection{Estudio del segundo mayor valor singualar de las matrices $\mathbf{D}_n$.}

\noindent $\bullet$ {\bf  ESTUDIO DEL COMPORTAMIENTO ASINTÓTICO de: }
$$
  \lim_{n\to \infty} \sigma_{2,n},
$$

\noindent tiene relación con el soporte de las medidas???

\noindent En los experimentos realizados la conjetura principal es:


\noindent {\bf CONJETURA PRINCIPAL:}


\noindent En un espacio de Sobolev con una derivada, es decir, del tipo
$$
  <p(t),q(t)>=\int p(t)q(t)\; d\mu_0 + \int p'(t)q'(t)d\mu_1
$$

\noindent asociada a medidas con soporte acotado y siendo
$$
  R=\max \{ \vert t \vert : t\in {\it supp}(\mu_0)\cup {\it supp}(\mu_1)\}
$$


\noindent se tiene que:

$$
  \lim_{n\to \infty} \sigma_{2,n} = R
$$


\noindent El estudio del segundo valor singular de una matriz ( y de cualquiera de ellos) tiene una expresión min max por el Teorema de Courant Fisher

\noindent A partir de estos teoremas se obtiene la siguiente caracterización de dicho valor singular como el mínimo  de las normas del operador restringido a subespacios vectoriales de dimensión $n-1$ o codimensión $1$, es decir,
$$
  \sigma_{2,n}^{2} = \min_{S: dim(S)=n-1} \max_{x\in S,\Vert x\Vert=1} <\mathbf{D}_n^t\mathbf{D}_n x,x>= \min_{S: dim(S)=n-1} \Vert \mathbf{D}_n/S\Vert^2.
$$







\noindent podría haber un min-max infinito dimensinal en el sentido de considerar el mínimo de las normas del operador multiplicación por $z$ entre todos los subespacios de codimensión $1$.

$$
  \min_{S: cod(S)=1} \Vert \mathcal{D}/S \Vert
$$

\noindent Puede ocurrir que

$$
  \min_{S: cod(S)=1} \Vert \mathcal{D}/S \Vert < \Vert \mathcal{D} \Vert = \infty.
$$

\noindent Sí, puede ocurrir, basta considerar un caso como una matriz que sea diagonal salvo en la primera fila que esté llena de unos, entonces las secciones tendrán normas que van a infinito pero si consideramos las restricciones en las sucesiones con $x_0=0$ se tendría que está acotada, y el mínimo anterior sería todavía menor.


$$
  \mathbf{D}= \begin{pmatrix}
    1 & 1           & 1           & \hdots & 1      & \hdots    \\
    1 & \frac{1}{2} & 0           & \hdots & 0      & \hdots
    \\
    0 & \frac{1}{2} & \frac{1}{3} & \vdots & \hdots & \hdots
    \\
    0 & 0           & 0           & \vdots & \hdots & \hdots  \
  \end{pmatrix}
$$


\noindent {\bf Problema:} ¿Es
$$
  \lim_{n\to \infty} \sigma_{2,n}= \min_{S: cod(S)=1} \Vert \mathcal{D}/S \Vert ????
$$
\subsection{Acotación de $\mathcal{D}$ en hiperplanos en espacios de  Sobolev con peso}

\noindent Un resultado muy interesante.

\begin{prop} Sea $w(t)$ un peso continuo en $[a,b]$ y consideremos el producto Sobolev con peso en $0$ dado por:
  $$
    \Vert p(t) \Vert^{2} = \int_{a}^{b} p^{2}(t)w(t)\; dt + p'(0)^{2}.
  $$



  \noindent Se tiene :
  \begin{enumerate}
    \item El operador multiplicación por $z$ no es acotado, por lo que:
          $$
            \lim_{n\to \infty} \sigma_{1,n} = \infty.
          $$
    \item  El operador multiplicación por $z$ es acotado en el subespacio infinito dimensional de codimensión $1$, $\mathbb{P}_{0}[z]$, y además
          $$
            \lim_{n\to \infty} \sigma_{2,n}= \min_{S: cod(S)=1} \Vert \mathcal{D}/S \Vert \leq \max\{\vert a \vert, \vert b \vert \}
          $$
  \end{enumerate}
\end{prop}

\begin{proof} Consideremos el subespacio vectorial de codimensión $1$, $\mathbb{P}_{0}[z]$, sea $p(z)\in \mathbb{P}[z]$ con $p(0)=0$, se tiene que:
  $$
    \Vert \mathcal{D}(p(t))\Vert^{2} = \int_{a}^{b} t^{2}p^{2}(t)\; dt + (tp(t))'(0)=  \int_{a}^{b} t^{2}p^{2}(t)\; dt + (p(t)+tp'(t))(0)\leq
  $$
  $$
    \max \{a^2,b^2\} \int_{a}^{b} p^{2}(t)\; dt + p(0)^2= \max \{a^2,b^2\} \int_{a}^{b} p^{2}(t)\; dt \leq
  $$
  $$
    \max \{a^2,b^2\} ( \int_{a}^{b} p^{2}(t)\; dt+p'(0)^2)=\max \{a^2,b^2\} \Vert p(t)\Vert^{2},
  $$

  \noindent por continuidad se tiene que $\mathcal{D}$ es acotado en $\mathbb{P}_0[z]$ siendo su norma menor o igual que $R=\max\{\vert a \vert, \vert b \vert \}$.
\end{proof}


\noindent {\bf Experimentos de Gabriel en este caso }:

$$
  \Vert p(t) \Vert^{2} = \int_{0}^{1} p^{2}(t)\; dt + p'(0)^{2}.
$$


\begin{enumerate}
  \item $\lim_{n \to \infty} \sigma_{1,n}=\infty$
  \item $$\lim_{n\to \infty} \sigma_{2,n}=1$$

        \noindent Hemos probado que en este caso:
        $$
          \lim_{n\to \infty} \sigma_{2,n} \leq 1
        $$
        \noindent INTENTAR PROBAR QUE ES $1$ !!!!

  \item
        $$
          \lim_{n\to \infty} \rho(\mathbf{D}_n)=1
        $$
\end{enumerate}


\noindent {\bf POSIBLES GENERALIZACIONES:}
\begin{prop} En un espacio de Sobolev con peso del tipo:
  $$
    \Vert p(z) \Vert^{2} = \int \vert p(z) \vert^{2} d\mu +  \vert p'(0) \vert^{2}
  $$

  \noindent tal que $0$ no es punto de evaluación de la medida $\mu$  se tiene que:

  \begin{enumerate}
    \item $\mathcal{D}$ no es acotado, por lo que:
          $$
            \lim_{n\to\infty} \sigma_{1,n}= \infty.
          $$
    \item $\mathcal{D}$ es acotado en algún hiperplano, por lo que:

          $$
            \lim_{n\to \infty} \sigma_{2,n}= \min_{S: cod(S)=1} \Vert \mathcal{D}/S \Vert \leq \max\{\vert z \vert : z\in {\it supp}(\mu) \}
          $$
  \end{enumerate}

\end{prop}

\subsection{Estudio detallado del caso Legendre+peso en 0}

\noindent Consideremos el siguiente producto de Sobolev:
$$
  \Vert p(t)\Vert^{2} = \int_{0}^{1} p^2(t)\; dt + p'(0)^{2}
$$


\noindent Nuestro objetivo es probar de forma sencilla que:
\begin{enumerate}
  \item $\mathcal{D}$ no es acotado y de hecho
        $$
          \lim_{n \to \infty} \sigma_{1,n}=\infty
        $$

  \item
        $$\min_{S: cod(S)=1} \Vert \mathcal{D}/S \Vert < \infty
        $$

  \item $$¿\lim_{n\to\infty} \rho(\mathbf{D}_n)= ???$$
\end{enumerate}


\section{Convergencia del  segundo valor singular máximo.}

\noindent Problemas:  Si tenemos una matriz hermitiana definida positiva, acotada, y consideramos las secciones, y sus valores singulares ordenados:
$$
  0 \leq \sigma_{n,n}\leq \dots \sigma_{2,n} \leq \sigma_{1,n}
$$

\begin{enumerate}
  \item ¿ es cierto que $\lim_{n\to \infty} \sigma_{2,n}$ existe???

        \noindent La respuesta es positiva y es consecuencia del Teorema de Cauchy de entralazamiento de los autovalores.

        \noindent {\it Cauchy interlace theorem} Let $A$ be an Hermitian matrix of order $n$ and let $B$ a principal submatrix of $A$ or order $n-1$. If $\lambda_n\leq \dots \leq \lambda_2\leq \lambda_1$ lists the eigenvalues of $A$ and $\mu_{n-1}\leq \dots \leq \mu_2\leq \mu_1$ the eigenvalues of $B$ then
        $$
          \lambda_n \leq \mu_{n-1} \leq \dots \leq \mu_2 \leq \lambda_1
        $$

        \noindent Este resultado justifica el entrelazamiento de las raíces de los polinomios ortogonales cuando consideramos medidas en el caso real puesto que las  matrices de Jacobi, que son las matrices de Hessemberg en estos casos, son hermitianas, y sus autovalores son las raíces de los polinomios ortogonales.

        \noindent Ahora, se puede probar que
        $\mathbf{D}^{t}_{n-1}\mathbf{D}_n$ es una submatriz de $\mathbf{D}_n^{t}\mathbf{D}_n$
        por lo que hay un entrelazamiento entre los valores singulares de $\mathbf{D}_{n-1}$ y los valores singulares de $\mathbf{D}_n$:
        $$
          \sigma_{n,n} \leq \sigma_{n-1,n-1} \leq \dots \leq \sigma_{2,n-1} \leq   \sigma_{2,n} \leq \sigma_{1,n-1} \leq \sigma_{1,n}.
        $$

        \noindent En particular, la sucesión $\{\sigma_{2,n}\}_{n=1}^{\infty}$ es una sucesión creciente, por lo que o converge o tiene límite infinito. EXISTE
        $$
          \lim_{n\to \infty} \sigma_{2,n} \leq \infty.
        $$
  \item ?????Coincide con
        $$
          \min_{V, cod(V)=1} \max \Vert \mathcal{D}/V \Vert
        $$
        \noindent

  \item
\end{enumerate}

\noindent Sería interesante probar (o buscar referencias) de que
$$
  \lim_{n\to \infty} \sigma_{2,n} = \inf_{V, cod(V)=1} \sup \Vert \mathcal{D}/V \Vert
$$



\newpage

\section{CASO DE CIRCUNFERENCIAS}

$$
  <p,q>=\int p(z) \overline{q(z)}d{\bf m} + \int_{\mathbf{S}_r} p'(z) \overline{q'(z)}d{\bf m_{0;r}}
$$

$$
  \Vert p(z) \Vert^{2}_{\mathbf{S}} = \Vert zp(z) \Vert_{{\bf m}}^{2} + \Vert p'(z) \Vert^{2}_{{\bf m}_{0;r}}=
$$






$$
  \Vert \mathcal{D}(p(z))\Vert_{\mathbf{S}}^{2} = \Vert \vert zp(z)\Vert_{{\bf m}}^{2} + \Vert (zp(z))' \Vert^{2}_{{\bf m}_{0;r}}=
$$

\noindent Puesto que en la primera integral $\vert z \vert=1$ se tiene que:

$$
  =\Vert  p(z) \Vert_{{\bf m}}
  ^{2} + \Vert p(z)+zp'(z) \Vert^{2}_{{\bf m}_{0;r}}
$$
$$
  \leq \Vert  p(z) \Vert^2_{{\bf m}} + 2 \Vert  p(z) \Vert^2_{{\bf m}_{0;r}} + 2\Vert zp'(z) \Vert^{2}_{{\bf m}_{0;r}} \leq \Vert  p(z) \Vert^2_{{\bf m}} + 2 \Vert  p(z) \Vert^2_{{\bf m}_{0;r}} + 2r^2\Vert p'(z) \Vert^{2}_{{\bf m}_{0;r}} \leq
$$
\noindent Si $C= \max \{1,2r^2\}$ se tiene que
$$
  \leq C(\Vert  p(z) \Vert^2_{{\bf m}}+\Vert p'(z) \Vert^{2}_{{\bf m}_{0;r}})+2\Vert  p(z) \Vert^2_{{\bf m}_{0;r}}= C\Vert p(z) \Vert^{2}_{\mathbf{S}}+ \Vert  p(z) \Vert^2_{{\bf m}_{0;r}}
$$

\noindent Utilizaremos ahora la desigualdad polinomial de Wirtinger para la medida ${\bf m}_{0;r}$ que dice que si $q(z)$ es un polinomio con $q(0)=0$ se tiene que:
$$
  \Vert q(z) \Vert^{2} = \int_{\mathbf{S}_r} \left(\sum_{k=0}^{n} a_kz^k \right)= \sum_{k=0}^{n} r^{2k} \vert a_k \vert^{2} \leq
$$
$$
  \sum_{k=0}^{n} r^{2k}  \vert a_k \vert^{2}\leq \sum_{k=1}^{n} kr^{2k-2}\vert a_k \vert^{2} = \int_{\mathbf{S}_r} \vert q'(z)\vert^{2}d{\bf m}_{0;r}=  \Vert q'(z)\Vert^{2}.
$$


\noindent Luego,
$$
  \vert q(0) \vert
$$


\end{document}
